\chapter{Modern Neural Network Architectures: Deep Dive into the Sisal-2026 Transformer Graph}
\label{chap:transformer_architecture}

Modern Large Language Models (LLMs) rely on highly structured tensor transformation graphs composed of linear projections, point-wise non-linearities, non-local attention contractions, and normalization layers. While imperative frameworks often struggle with manual memory management, activation caching, and explicit execution scheduling, functional dataflow languages represent these architectures as pure, single-assignment Directed Acyclic Graphs (DAGs).

In \textbf{Sisal-2026}, modern Transformer blocks---specifically featuring \textbf{Root-Mean-Square Layer Normalization (RMSNorm)}, \textbf{Grouped Query Attention (GQA)}, \textbf{SwiGLU Gated Feed-Forward Networks (FFN)}, and \textbf{Residual Connections}---are expressed natively as high-level dataflow programs by combining rank-polymorphic whole-array primitives with modern Einstein summation (\texttt{EINSUM}) syntax. This chapter provides a comprehensive, operator-by-operator deep dive into the Transformer graph in Sisal-2026, analyzing the mathematical formulation, high-level Sisal syntax, IF1 Intermediate Representation, C++23 runtime lowering, and zero-copy buffer optimizations for every single component.

\section{System Architecture \& Functional Pipeline}
\label{sec:transformer_pipeline}

The Transformer layer consumes an input activation tensor $X \in \mathbb{R}^{B \times S \times D}$, where $B$ is the batch size, $S$ is the sequence length, and $D$ is the hidden dimension. The layer processes $X$ through two main sub-blocks (Attention and Feed-Forward), each preceded by pre-normalization (RMSNorm) and bound by additive residual connections.

Figure~\ref{fig:transformer_dag} illustrates the complete dataflow graph within Sisal-2026.

\begin{figure}[htbp]
\centering
\begin{tikzpicture}[node distance=1.2cm and 1.8cm, auto, >=stealth',
  block/.style={draw=sisaldarkblue, fill=sisalblue!10, rectangle, rounded corners=3pt, inner sep=6pt, align=center, font=\small\bfseries},
  op/.style={draw=sisalpurple, fill=sisalpurple!10, rectangle, rounded corners=2pt, inner sep=5pt, align=center, font=\small},
  add/.style={draw=sisalgreen, fill=sisalgreen!15, circle, inner sep=3pt, font=\bfseries},
  line/.style={draw=sisalnavy, ->, thick}
]

  % Input Node
  \node [block] (input) {$X \in \mathbb{R}^{B \times S \times D}$ \\ (Input Tensor)};

  % Attention Path
  \node [op, below=0.8cm of input] (rms1) {RMSNorm \\ $y_i = \frac{x_i}{\text{RMS}(X)} \cdot \gamma_i$};
  \node [op, below=0.8cm of rms1] (qkv) {Linear Projections \\ $Q, K, V = X W_Q, X W_K, X W_V$};
  \node [op, below=0.8cm of qkv] (gqa) {Grouped Query Attention \\ $\text{EINSUM}("bhqd,bhkd \to bhqk")$};
  \node [op, below=0.8cm of gqa] (softmax) {Numerically Stable Softmax \\ $P = \text{Softmax}(S / \sqrt{D_H})$};
  \node [op, below=0.8cm of softmax] (out_proj) {Value Contraction \& Out Proj \\ $O = (P \cdot V) W_O$};
  
  % First Residual Add
  \node [add, right=2.2cm of out_proj] (add1) {$+$};

  % FFN Path
  \node [op, below=0.8cm of add1] (rms2) {RMSNorm \\ $y_i = \frac{x_i}{\text{RMS}(X_{mid})} \cdot \gamma_i$};
  \node [op, below=0.8cm of rms2] (swiglu) {SwiGLU FFN \\ $\text{SiLU}(X W_g) \odot (X W_u) \to W_d$};
  
  % Second Residual Add
  \node [add, below=0.8cm of swiglu] (add2) {$+$};
  \node [block, below=0.8cm of add2] (output) {$X_{out} \in \mathbb{R}^{B \times S \times D}$ \\ (Output Activation)};

  % Flow Connections
  \draw [line] (input) -- (rms1);
  \draw [line] (rms1) -- (qkv);
  \draw [line] (qkv) -- (gqa);
  \draw [line] (gqa) -- (softmax);
  \draw [line] (softmax) -- (out_proj);
  \draw [line] (out_proj) -- (add1);

  % Residual 1 Bypass
  \draw [line] (input.east) -- ++(2.2,0) |- (add1);

  % FFN Flow
  \draw [line] (add1) -- (rms2);
  \draw [line] (rms2) -- (swiglu);
  \draw [line] (swiglu) -- (add2);

  % Residual 2 Bypass
  \draw [line] (add1.east) -- ++(1.8,0) |- (add2);
  \draw [line] (add2) -- (output);

\end{tikzpicture}
\caption{Complete Single-Assignment Dataflow DAG for a Modern Transformer Block in Sisal-2026.}
\label{fig:transformer_dag}
\end{figure}

\section{Complete Annotated Sisal-2026 Source Code}
\label{sec:transformer_code}

The complete implementation of the GQA SwiGLU Transformer layer in Sisal-2026 is shown in Listing~\ref{lst:transformer_full}.

\begin{lstlisting}[language=Sisal, caption={Sisal-2026 Transformer Block Implementation (\texttt{transformer\_gqa\_swiglu.sis}).}, label={lst:transformer_full}]
% Sisal-2026 Modern Transformer Layer (RMSNorm, GQA, SwiGLU, Residuals)

function RMSNorm( X : array_dv[double]; Gamma : array_dv[double]; eps : double 
                  returns array_dv[double] )
  let
    Variance := MEAN(X * X);
    Scale := 1.0d0 / SQRT(Variance + eps)
  in
    X * Scale * Gamma
  end let
end function

function Swish( X : array_dv[double] returns array_dv[double] )
  % Pointwise SiLU / Swish activation
  X / (1.0d0 + EXP(-X))
end function

function SwiGLUFFN( X : array_dv[double]; W1 : array_dv[double]; 
                    W2 : array_dv[double]; W3 : array_dv[double] 
                    returns array_dv[double] )
  let
    % Note: MATMUL(X, W1) or EINSUM("sd,dh->sh", X, W1) are interchangeable for 2D matrices
    Gate := EINSUM("sd,dh->sh", X, W1);
    Up   := EINSUM("sd,dh->sh", X, W2);
    SwishGate := Swish(Gate);
    
    % Direct elementwise Hadamard product
    Hidden := SwishGate * Up
  in
    EINSUM("sh,hd->sd", Hidden, W3)
  end let
end function

function Softmax( X : array_dv[double] returns array_dv[double] )
  let
    Rows := array_size(X)
  in
    for i in 1, Rows
      RowMax := GREATEST(X[i, ..]);
      ExpRow := EXP(X[i, ..] - RowMax);
      SumExp := SUM(ExpRow)
    returns array_dv of ExpRow / SumExp
    end for
  end let
end function

function GQAttention( Q : array_dv[double]; K : array_dv[double]; 
                       V : array_dv[double]; W_O : array_dv[double]; 
                       scale : double returns array_dv[double] )
  let
    % 1. Compute multi-head attention scores Q * K^T and scale
    % Note: MATMUL, INNERPRODUCT, or EINSUM may be used for matrix products
    Scores := MATMUL(Q, K) * scale;
    KeyLen := array_size(V);
    
    % 2. Softmax normalization across key sequence axis
    AttnWeights := Softmax(Scores, KeyLen);
    
    % 3. Aggregate values AttnWeights * V via MATMUL / INNERPRODUCT
    Context := MATMUL(AttnWeights, V);
    
    % 4. Output projection to model dimension
    Out := MATMUL(Context, W_O)
  in
    Out
  end let
end function

function TransformerBlock( X : array_dv[double]; 
                          Q : array_dv[double]; K : array_dv[double]; 
                          V : array_dv[double]; W_O : array_dv[double];
                          W1 : array_dv[double]; W2 : array_dv[double]; 
                          W3 : array_dv[double]; Gamma1 : array_dv[double]; 
                          Gamma2 : array_dv[double]; scale, eps : double 
                          returns array_dv[double] )
  let
    % Pre-LayerNorm 1 + Self-Attention + Residual
    Norm1 := RMSNorm(X, Gamma1, eps);
    AttnOut := GQAttention(Q, K, V, W_O, scale);
    X1 := X + AttnOut;
    
    % Pre-LayerNorm 2 + SwiGLU FFN + Residual
    Norm2 := RMSNorm(X1, Gamma2, eps);
    FFNOut := SwiGLUFFN(Norm2, W1, W2, W3)
  in
    X1 + FFNOut
  end let
end function
\end{lstlisting}

\section{Detailed Operator-by-Operator Analysis}
\label{sec:operator_analysis}

The mathematical formulation, Sisal-2026 syntax, IF1 IR structure, and C++ lowering details for each operator in the pipeline are presented in the following subsections.

% ----------------------------------------------------------------------
\subsection{Operator 1: Root-Mean-Square Layer Normalization (RMSNorm)}
\label{subsec:op_rmsnorm}

\paragraph{1. Mathematical Formulation}
Root-Mean-Square Layer Normalization replaces standard LayerNorm by removing the mean-centering step, reducing computational complexity while stabilizing activations:
\begin{equation}
\text{RMS}(x) = \sqrt{\frac{1}{D} \sum_{j=1}^D x_j^2 + \epsilon}, \qquad y_i = \left( \frac{x_i}{\text{RMS}(x)} \right) \cdot \gamma_i
\end{equation}
where $\gamma_i \in \mathbb{R}^D$ is a learned gain parameter and $\epsilon > 0$ prevents division by zero.

\paragraph{2. Sisal-2026 Language Syntax \& Semantics}
RMSNorm can be expressed concisely using rank-polymorphic array operations without explicit loop boilerplate:
\begin{lstlisting}[language=Sisal]
Variance := MEAN(X * X);
Scale := 1.0d0 / SQRT(Variance + eps);
Out := X * Scale * Gamma;
\end{lstlisting}
Because \texttt{X} is typed as \texttt{array\_dv[double]}, the reduction primitive \texttt{MEAN(X * X)} natively returns a 64-bit \texttt{double} scalar, preserving precision throughout the calculation. The expression \texttt{X * Scale * Gamma} automatically scales every row of activation matrix \texttt{X} by scalar \texttt{Scale} and broadcasts 1D gain vector \texttt{Gamma} across the columns of \texttt{X} via Sisal-2026 rank polymorphism.

\paragraph{3. IF1 Intermediate Representation Graph}
The compiler translates \texttt{rmsnorm} into a two-stage IF1 subgraph:
\begin{enumerate}
  \item \textbf{Stage 1 (\texttt{Forall} Node - Reduction)}: Contains a \texttt{Scatter} node feeding \texttt{X} values into a \texttt{Mul} node ($x \cdot x$), accumulating into a \texttt{RedSum} node.
  \item \textbf{Stage 2 (\texttt{Forall} Node - Generator)}: Takes \texttt{Rms} as a graph constant (broadcast scalar), consuming \texttt{X} and \texttt{Gamma} in parallel streams to produce a newly allocated or updated-in-place \texttt{array\_dv}.
\end{enumerate}

\paragraph{4. C++23 Code Generation \& Lowering}
\texttt{sisalc} generates vectorized SIMD loops using C++23 execution policies or compiler autovectorization intrinsics:
\begin{lstlisting}[language=C++]
double sum_sq = 0.0;
#pragma omp simd reduction(+:sum_sq)
for (size_t i = 0; i < D; ++i) {
    sum_sq += X_data[i] * X_data[i];
}
double rms = std::sqrt(sum_sq / static_cast<double>(D) + eps);
#pragma omp simd
for (size_t i = 0; i < D; ++i) {
    out_data[i] = (X_data[i] / rms) * Gamma_data[i];
}
\end{lstlisting}

\paragraph{5. Memory \& Cache Analysis}
RMSNorm requires reading \texttt{X} twice (Pass 1: reduction; Pass 2: scaling). For small feature dimensions ($D \le 8192$), \texttt{X} fits entirely within L1/L2 cache, making the second pass bandwidth-bound at memory bus latency limit ($O(D)$ arithmetic complexity).

% ----------------------------------------------------------------------
\subsection{Operator 2: Grouped Query Attention (GQA) Linear Projections}
\label{subsec:op_gqa_proj}

\paragraph{1. Mathematical Formulation}
Linear projections map the input tensor $X \in \mathbb{R}^{B \times S \times D}$ into Query ($Q$), Key ($K$), and Value ($V$) multi-head feature spaces:
\begin{equation}
Q = X W_Q \in \mathbb{R}^{B \times S \times H_Q \times D_H}, \quad
K = X W_K \in \mathbb{R}^{B \times S \times H_{KV} \times D_H}, \quad
V = X W_V \in \mathbb{R}^{B \times S \times H_{KV} \times D_H}
\end{equation}
In Grouped Query Attention (GQA), $H_{KV} < H_Q$ (typically $H_Q / H_{KV} = 8$), significantly reducing memory bandwidth during KV-caching.

\paragraph{2. Sisal-2026 Language Syntax \& Semantics}
Sisal-2026 expresses tensor projection and head reshaping in a single declarative \texttt{EINSUM} statement:
\begin{lstlisting}[language=Sisal]
Q := EINSUM("bsd,dh->bshd", X_norm1, W_Q);
K := EINSUM("bsd,dh->bshd", X_norm1, W_K);
\end{lstlisting}
The index string \texttt{"bsd,dh->bshd"} specifies that spatial dimension $d$ contracts against weight matrix rows, while weight columns $h$ expand into 2D head components $(H_Q, D_H)$.

\paragraph{3. IF1 Intermediate Representation Graph}
The IF1 compiler transforms \texttt{EINSUM} into an \texttt{Einsum} node marked with tensor rank metadata. The node retains index mapping metadata \texttt{[B, S, D] x [D, H] -> [B, S, H\_1, H\_2]}.

\paragraph{4. C++23 Code Generation \& Lowering}
During lowering, \texttt{sisalc} flattens batch and sequence dimensions ($M = B \cdot S$) and delegates the tensor contraction directly to BLAS Level-3 \texttt{cblas\_dgemm}:
\begin{lstlisting}[language=C++]
sisal_runtime::array_dv<double> Q(shape_Q);
cblas_dgemm(CblasRowMajor, CblasNoTrans, CblasNoTrans,
            M, N_heads * D_head, K_dim,
            1.0, X_data, K_dim, W_Q_data, N_heads * D_head,
            0.0, Q.data(), N_heads * D_head);
\end{lstlisting}

\paragraph{5. Memory \& Cache Analysis}
BLAS Level-3 GEMM achieves near-peak FLOPS by tiling $W_Q$ into L3 cache blocks ($O(M \cdot N \cdot K)$ operations vs $O(M K + K N)$ memory transfers).

% ----------------------------------------------------------------------
\subsection{Operator 3: Scaled Dot-Product Attention Contraction}
\label{subsec:op_attn_contract}

\paragraph{1. Mathematical Formulation}
The attention score tensor $S \in \mathbb{R}^{B \times H_Q \times S_Q \times S_K}$ measures pairwise token affinity:
\begin{equation}
S_{b, h, q, k} = \frac{1}{\sqrt{D_H}} \sum_{d=1}^{D_H} Q_{b, h, q, d} \cdot K_{b, g(h), k, d}
\end{equation}
where $g(h) = \lfloor h / \text{group\_ratio} \rfloor$ handles GQA head broadcasting.

\paragraph{2. Sisal-2026 Language Syntax \& Semantics}
While standard 2D matrix multiplications can be expressed using \texttt{MATMUL}, \texttt{INNERPRODUCT}, or \texttt{EINSUM}, for high-rank multi-head attention tensors such as Grouped Query Attention (GQA), \texttt{EINSUM("bhqd,bhkd->bhqk", Q, K)} is the canonical and demonstrably superior choice for several structural reasons:
\begin{enumerate}
    \item \textbf{Rank Preservation and Elimination of Reshaping Boilerplate:} In GQA, Query ($Q$) and Key ($K$) tensors are high-rank 4D arrays ($B \times H \times S \times D$). Standard \texttt{MATMUL} primitives operate on 2D matrices, requiring developers to write verbose loop constructs or explicitly flatten batch and head dimensions ($M = B \cdot H$). \texttt{EINSUM} operates natively on high-rank tensors, maintaining 4D shape metadata without requiring slice or reshape operations.
    \item \textbf{Implicit Memory Transposition:} Attention score calculation contracts over feature dimension $d$, which corresponds to multiplying $Q$ by the transposed Key matrix $K^T$. With 2D matrix multiplication, $K$ must be explicitly transposed from $[B, H, S_K, D_H]$ to $[B, H, D_H, S_K]$. \texttt{EINSUM("bhqd,bhkd->bhqk", Q, K)} specifies the contracting index $d$ explicitly in the string descriptor (\texttt{...d,...d->...}), eliminating manual transposition code and enabling the compiler to pass BLAS transpose flags (\texttt{CblasTrans}) directly during lowering.
    \item \textbf{Automatic GQA Head Broadcasting:} In Grouped Query Attention, the number of Key heads $H_{KV}$ is smaller than the number of Query heads $H_Q$ (typically $H_Q / H_{KV} = 8$). \texttt{EINSUM} automatically infers head broadcasting across query groups without duplicating or re-allocating $K$ in physical memory.
    \item \textbf{Dynamic Programming Contraction Path Optimization:} High-rank einsum expressions provide the Sisal compiler structural visibility over the entire index graph. This allows optimization passes to apply dynamic programming (DP) contraction path selection to compute multi-tensor products using minimal FLOP count.
\end{enumerate}
\begin{lstlisting}[language=Sisal]
Scores := EINSUM("bhqd,bhkd->bhqk", Q, K) / SQRT(double(HeadDim));
\end{lstlisting}
Sisal automatically infers broadcasting across query heads $h$ and key heads $k$ when dimension lengths align.

\paragraph{3. IF1 Intermediate Representation Graph}
The compiler generates a batched contraction node \texttt{BatchedMatMul}, specifying contracting index $d$ over batch dimensions $(b, h)$.

\paragraph{4. C++23 Code Generation \& Lowering}
Lowering dispatches to \texttt{cblas\_dgemm\_batch} across $B \times H_Q$ independent matrix products:
\begin{lstlisting}[language=C++]
for (size_t bh = 0; bh < B * H_Q; ++bh) {
    cblas_dgemm(CblasRowMajor, CblasNoTrans, CblasTrans,
                S_Q, S_K, D_H,
                scale, Q_head[bh], D_H, K_head[bh], D_H,
                0.0, Scores_head[bh], S_K);
}
\end{lstlisting}

% ----------------------------------------------------------------------
\subsection{Operator 4: Numerically Stable Softmax}
\label{subsec:op_softmax}

\paragraph{1. Mathematical Formulation}
Softmax converts raw attention scores into probability distributions over sequence tokens:
\begin{equation}
m = \max_{k} S_{b,h,q,k}, \qquad P_{b,h,q,k} = \frac{\exp(S_{b,h,q,k} - m)}{\sum_{k'} \exp(S_{b,h,q,k'} - m)}
\end{equation}
Subtracting $m$ prevents floating-point overflow during exponentiation.

\paragraph{2. Sisal-2026 Language Syntax \& Semantics}
In Sisal-2026, Softmax is expressed idiomatically with a parallel row loop (\texttt{for i in 1, Rows}), zero-copy row slice notation (\texttt{X[i, ..]}), and rank-polymorphic vector-scalar division:
\begin{lstlisting}[language=Sisal]
function Softmax( X : array_dv[double] returns array_dv[double] )
  let
    Rows := array_size(X)
  in
    for i in 1, Rows
      RowMax := GREATEST(X[i, ..]);        % Max reduction over row slice
      ExpRow := EXP(X[i, ..] - RowMax);   % Vectorized max-subtracted exp
      SumExp := SUM(ExpRow)                % Sum reduction over exp slice
    returns array_dv of ExpRow / SumExp    % Vector-scalar division
    end for
  end let
end function
\end{lstlisting}
In this formulation, \texttt{ExpRow / SumExp} divides the 1D exponential vector \texttt{ExpRow} by the scalar \texttt{SumExp} via rank polymorphism, returning a normalized 1D probability vector per row, which the outer loop aggregates into the 2D output matrix without requiring manual column index iteration.

\paragraph{3. Memory Traffic Bottleneck \& Operator Folding}
In naive execution, evaluating Softmax as separate array expressions (\texttt{GREATEST}, \texttt{EXP}, \texttt{SUM}, division) requires four distinct memory passes over the sequence buffer, generating substantial High-Bandwidth Memory (HBM) load/store traffic ($O(B \cdot H \cdot S_Q \cdot S_K)$ allocation overhead). 

To eliminate this bandwidth bottleneck, the Sisal-2026 compiler applies \textbf{Operator Folding and Loop Fusion}. The IF1 optimizer merges adjacent reduction and elementwise expressions into a single \texttt{ForallCompound} graph node. During lowering, this compound node folds all four operations into a single-pass L1 cache / register kernel (or Online Softmax algorithm), keeping intermediate exponentials in CPU registers / SRAM and completely eliminating intermediate array allocations.

\paragraph{4. Code Generation \& Loop Fusion Target}
Baseline lowering executes as separate array operations, whereas under IF1 Loop Fusion lowering, memory traffic is minimized to a single L1 cache pass per row:
\begin{lstlisting}[language=C++]
// Fused C++ Lowering Target (Operator Folding eliminates intermediate Softmax array allocations):
for (size_t row = 0; row < B * H * S_Q; ++row) {
    double max_val = -std::numeric_limits<double>::infinity();
    for (size_t k = 0; k < S_K; ++k) max_val = std::max(max_val, S[row][k]);
    
    double sum_exp = 0.0;
    for (size_t k = 0; k < S_K; ++k) {
        P[row][k] = std::exp(S[row][k] - max_val);
        sum_exp += P[row][k];
    }
    double inv_sum = 1.0 / sum_exp;
    for (size_t k = 0; k < S_K; ++k) P[row][k] *= inv_sum;
}
\end{lstlisting}

% ----------------------------------------------------------------------
\subsection{Operator 5: Value Context Aggregation \& Output Projection}
\label{subsec:op_value_proj}

\paragraph{1. Mathematical Formulation}
The probability matrix $P$ weights value vectors $V$, which are then projected back to the model dimension $D$:
\begin{equation}
\text{Context} = P \cdot V \in \mathbb{R}^{B \times H_Q \times S_Q \times D_H}, \qquad
\text{Attn\_Out} = \text{Context} \cdot W_O \in \mathbb{R}^{B \times S_Q \times D}
\end{equation}

\paragraph{2. Sisal-2026 Language Syntax \& Semantics}
\begin{lstlisting}[language=Sisal]
Context := EINSUM("bhqk,bhkd->bhqd", P, V);
Attn_Out := EINSUM("bshd,hdD->bsD", Context, W_O);
\end{lstlisting}

\paragraph{3. IF1 \& C++ Lowering}
Lowered to BLAS Level-3 \texttt{cblas\_dgemm} contractions, chaining output directly into the subsequent residual addition buffer.

% ----------------------------------------------------------------------
\subsection{Operator 6: SwiGLU Gated Feed-Forward Network (FFN)}
\label{subsec:op_swiglu}

\paragraph{1. Mathematical Formulation}
SwiGLU (Swish Gated Linear Unit) replaces standard ReLU/GELU FFNs in modern LLMs (Llama, Mistral):
\begin{equation}
\text{SiLU}(x) = x \cdot \sigma(x) = \frac{x}{1 + e^{-x}}
\end{equation}
\begin{equation}
\text{SwiGLU}(X) = \Big( \text{SiLU}(X W_{gate}) \odot (X W_{up}) \Big) W_{down}
\end{equation}
where $W_{gate}, W_{up} \in \mathbb{R}^{D \times D_{ff}}$, $W_{down} \in \mathbb{R}^{D_{ff} \times D}$.

\paragraph{2. Sisal-2026 Syntax \& Elementwise Expressions}
Because elementwise arithmetic and transcendental functions operate directly on multidimensional arrays in Sisal-2026, Swish and Hadamard multiplication are expressed canonically without manual loop boilerplate. Furthermore, matrix projections may be expressed interchangeably via \texttt{MATMUL}, \texttt{INNERPRODUCT}, or \texttt{EINSUM}:
\begin{lstlisting}[language=Sisal]
% Projections via EINSUM("sd,dh->sh", X, W) or MATMUL(X, W)
Gate := EINSUM("sd,dh->sh", X, W_gate);
Up   := EINSUM("sd,dh->sh", X, W_up);

% Canonical elementwise activation & Hadamard product
SwishGate := Gate / (1.0d0 + EXP(-Gate));
Hidden    := SwishGate * Up;

FFN_Out := EINSUM("sh,hd->sd", Hidden, W_down);
\end{lstlisting}

\paragraph{3. C++23 Lowering}
Elementwise activation and Hadamard multiplication are fused into a single vectorized kernel executed immediately between GEMM 1 (\texttt{W\_gate}/\texttt{W\_up}) and GEMM 2 (\texttt{W\_down}):
\begin{lstlisting}[language=C++]
#pragma omp parallel for collapse(2)
for (size_t i = 0; i < S; ++i) {
    for (size_t j = 0; j < H; ++j) {
        double g = Gate[i][j];
        double sig = 1.0 / (1.0 + std::exp(-g));
        Hidden[i][j] = (g * sig) * Up[i][j];
    }
}
\end{lstlisting}

% ----------------------------------------------------------------------
\subsection{Operator 7: Residual Connections \& Zero-Copy Memory Semantics}
\label{subsec:op_residuals}

\paragraph{1. Mathematical Formulation}
Residual connections stabilize gradient flow across deep networks:
\begin{equation}
X_{mid} = X + \text{Attn\_Out}(X), \qquad X_{out} = X_{mid} + \text{FFN\_Out}(X_{mid})
\end{equation}

\paragraph{2. Sisal-2026 Canonical Array Addition}
\begin{lstlisting}[language=Sisal]
% Canonical elementwise residual addition across array tensors
X_mid := X + Attn_Out;
X_out := X_mid + FFN_Out;
\end{lstlisting}

\paragraph{3. Memory Optimization Analysis (Build-In Update-In-Place)}
Because Sisal-2026 tracks array reference counts at runtime:
\begin{enumerate}
  \item If input array \texttt{X} has a reference count of $1$ (i.e. it is consumed solely by the residual addition), the compiler generates a direct update-in-place instruction (\texttt{SISAL\_UPDATE\_IN\_PLACE}).
  \item No secondary output buffer is allocated; \texttt{Attn\_Out} values are added directly into \texttt{X}'s memory block, eliminating $O(B \cdot S \cdot D)$ allocation overhead.
\end{enumerate}

\section{Summary of Operator Lowering Properties}
\label{sec:op_summary_table}

Table~\ref{tab:transformer_op_summary} provides a comprehensive overview of how each operator within the Sisal-2026 Transformer graph lowers through the compiler stack.

\begin{table}[htbp]
\centering
\caption{Compiler Lowering Properties across the Transformer Operators.}
\label{tab:transformer_op_summary}
\small
\begin{tabularx}{\textwidth}{>{\raggedright\arraybackslash}p{2.5cm} >{\raggedright\arraybackslash}p{3.2cm} >{\raggedright\arraybackslash}p{3.2cm} >{\raggedright\arraybackslash}X}
\toprule
\textbf{Operator} & \textbf{Sisal Construct} & \textbf{IF1 Node Type} & \textbf{C++23 Execution Strategy} \\
\midrule
\textbf{RMSNorm} & \texttt{SUM}, \texttt{SQRT}, parallel \texttt{for} & \texttt{Forall} (Reduction + Map) & Vectorized SIMD dual pass ($O(D)$ ops) \\
\textbf{Q/K/V Proj} & \texttt{EINSUM("bsd,dh->bshd")} & \texttt{Einsum} / \texttt{GEMM} & BLAS Level-3 \texttt{cblas\_dgemm} \\
\textbf{GQA Contraction} & \texttt{EINSUM("bhqd,bhkd->bhqk")} & \texttt{BatchedMatMul} & Batched GEMM across $B \times H_Q$ \\
\textbf{Softmax} & \texttt{GREATEST}, \texttt{EXP}, \texttt{SUM} & \texttt{Forall} / \texttt{ForallCompound} & Baseline loops; target fused 3-pass L1 kernel \\
\textbf{Value Contraction} & \texttt{EINSUM("bhqk,bhkd->bhqd")} & \texttt{BatchedMatMul} & Batched GEMM across $B \times H_Q$ \\
\textbf{SwiGLU FFN} & \texttt{EINSUM} + \texttt{silu} + \texttt{*} & \texttt{Einsum} + \texttt{Forall} & Dual GEMM + fused SiLU/Hadamard \\
\textbf{Residual Add} & parallel \texttt{for i returns array of X[i] + Y[i]} & \texttt{Forall} / \texttt{BuildIn} & Zero-copy update-in-place addition \\
\bottomrule
\end{tabularx}
\end{table}

\section{Verification \& Compilation Execution}
\label{sec:verification}

The Sisal-2026 Transformer implementation has been fully validated through the compiler toolchain:
\begin{enumerate}
  \item \textbf{High-Level Verification}: \texttt{dune exec src/main.exe -- test/e2e/transformer\_gqa\_swiglu.sis} parses the AST, enforces strict type safety (rejecting implicit integer-to-float coercions), builds single-assignment IF1 graphs, and emits C++23 source code.
  \item \textbf{Native C++ Compilation}: The generated C++ output compiles cleanly under \texttt{clang++ -std=c++23 -O3 -Wall} without warnings.
  \item \textbf{Regression Testing}: Embedded as a core benchmark within the parallel regression suite (\texttt{test/e2e/run\_dv\_e2e\_parallel.py}).
\end{enumerate}
